\section{Chapitre~\ref{sec:acet}. Ajoutons \nt{ACET}}

Les trois actions de l'acétylcholine sont mesurées sur tranches et dans le cerveau vivant, le conflit entre écriture et lecture est montré sur le modèle du livre, et que \nt{ACET} serve de ligne de validation d'écriture et signale l'incertitude attendue est une hypothèse en partie appuyée par l'expérience.

{\footnotesize
\setlength{\tabcolsep}{3pt}\renewcommand{\arraystretch}{1.15}
\begin{longtable}{>{\raggedright}p{0.5cm}>{\raggedright}p{4.0cm}>{\raggedright}p{5.6cm}>{\raggedright}p{2.3cm}>{\raggedright\arraybackslash}p{1.7cm}}
\toprule
N\textsuperscript{o} & Affirmation & Expérience et ce qui est mesuré & Source & Statut \\
\midrule\endfirsthead
\toprule
N\textsuperscript{o} & Affirmation & Expérience et ce qui est mesuré & Source & Statut \\
\midrule\endhead
1 & Les neurones \nt{ACET} du cerveau sont peu nombreux, réunis en quelques groupes, et leurs sorties se répartissent dans tout le cortex: c'est un canal de diffusion & Marquage par anticorps de l'enzyme de synthèse de l'acétylcholine et traçage rétrograde chez le rat: le cortex, l'hippocampe, l'amygdale, le bulbe olfactif et le thalamus reçoivent l'acétylcholine surtout de six groupes compacts de cellules du télencéphale basal et du tronc cérébral & \cite{Mesulam1983} & mesuré \\
2 & Le niveau d'acétylcholine dans le cortex est élevé à l'éveil et bas en sommeil lent & Microdialyse dans le cortex et l'hippocampe de chats libres de leurs mouvements, avec enregistrement EEG: la libération d'acétylcholine est supérieure de $\sim100\%$ à l'éveil calme et de $\sim175\%$ à l'éveil actif par rapport au sommeil lent; en sommeil paradoxal, dans le cortex, elle est la même qu'à l'éveil & \cite{Marrosu1995,Hasselmo1999} & mesuré \\
3 & Le sens du signal est fixé par le récepteur: l'acétylcholine a des récepteurs-canaux rapides et des récepteurs lents qui déclenchent des réactions dans la cellule (proposition~\ref{prop:receptor}) & Revue de mesures: l'action de l'acétylcholine sur l'excitabilité, la transmission et la plasticité dépend du lieu de libération, du sous-type de récepteur (nicotiniques: canaux ioniques; muscariniques: par des protéines G) et de la cellule destinataire & \cite{Picciotto2012} & mesuré \\
4 & Première action: affaiblir les connexions internes sans presque toucher aux entrées extérieures (facteur $1-a$ dans~\eqref{eq:acetnet}) & Tranches de cortex olfactif de rat: à $100$\,\textmu M d'agonistes de l'acétylcholine, les potentiels des connexions internes (couche Ib) baissent de plus de $60\%$, ceux de l'entrée externe (couche Ia) de moins de $15\%$, dans la même cellule; la suppression est présynaptique. Tranches d'hippocampe (CA1): $89\%$ contre $40\%$ pour la voie interne et la voie d'entrée & \cite{HasselmoBower1992,HasselmoSchnell1994} & mesuré \\
5 & Deuxième action: renforcer l'entrée (facteur $\frac{1+a}{2}$) & Tranches de néocortex: les récepteurs nicotiniques renforcent les synapses de l'entrée thalamique et n'agissent pas sur les synapses intracorticales; les muscariniques affaiblissent les deux types. L'acétylcholine augmente l'excitabilité des neurones (revue) & \cite{Gil1997,Hasselmo2006} & mesuré \\
6 & Troisième action: faciliter la modification des poids (vitesse $\eta a$) & Rat: la stimulation électrique du noyau basal (source de l'acétylcholine du cortex), appariée à un son, provoque chez l'animal adulte une forte réorganisation de la carte du cortex auditif en faveur du son présenté & \cite{KilgardMerzenich1998,Hasselmo2006} & mesuré \\
7 & La règle de Hebb pour motifs clairsemés, dans la seconde équation~\eqref{eq:acetnet}, donne une mémoire associative qui complète un motif à partir d'une partie & Calcul de la capacité d'un réseau à motifs clairsemés et règle de covariance: pour une faible fraction de neurones actifs $p$, le réseau stocke nettement plus de motifs que pour $p=1/2$ & \cite{Tsodyks1988} & théorie \\
8 & Écrire avec $a$ bas abîme la mémoire: le réseau écrit ce dont il se souvient, et non ce qu'il voit; pour $a=0.1$, les dégâts ne sont pas moindres que pour $a=0$ & \texttt{models/\allowbreak acet\_memory.py}: $200$ neurones, cinq motifs de $20$, un nouveau motif partage $10$ neurones avec l'un d'eux, $10$ simulations. Pour $a=0$, le nouveau motif n'est pas mémorisé, l'ancien entraîne $5.9$ neurones étrangers sur $10$; pour $a=0.1$, le nouveau motif est rappelé ($0.97$), mais les deux fusionnent: le nouveau entraîne $7.3$ neurones étrangers, l'ancien $7.9$; à partir de $a=0.2$, moins d'un & démonstration dans le chapitre & modèle \\
9 & Proposition~\ref{prop:writeread}: aucun niveau $a$ ne convient aussi bien à l'écriture qu'à la lecture (fig.~\ref{fig:acetsim}) & Même script, vitesse d'apprentissage $\eta a$: le nouveau motif est entièrement mémorisé en une présentation pour $a\ge0.9$, à $0.8$ pour $a=0.75$, pas du tout pour $a\le0.5$. Lecture d'après la moitié du motif: $1.0$ pour $a\le0.2$, $0.96$ pour $a=0.3$, $0.04$ pour $a=0.4$ & démonstration dans le chapitre & modèle \\
10 & Dans le cerveau, un seul fil de diffusion, \nt{ACET}, commute écriture et lecture & L'idée vient de modèles construits à partir de mesures sur tranches. L'expérience la plus directe est humaine: la scopolamine, qui bloque les récepteurs muscariniques, gêne l'apprentissage de nouvelles paires de mots, surtout celles qui recoupent d'anciennes, mais pas le rappel des paires déjà apprises. Les deux modes du réseau n'ont pas été mesurés directement & \cite{HasselmoAndersonBower1992,Atri2004} & hypothèse \\
11 & Le filtre~\eqref{eq:kalman} est optimal; le poids de l'entrée et la vitesse d'apprentissage sont une même grandeur $P/R$; en régime établi, $P=\sqrt{qR}$ et la mémoire vaut $\sqrt{R/q}$ (proposition~\ref{prop:kalman}) & Aucune expérience n'est requise: l'optimalité du filtre de Kalman--Bucy est un résultat classique de la théorie de l'estimation; le régime établi est démontré dans le chapitre & \cite{KalmanBucy1961} & théorie \\
12 & \nt{ACET} communique au réseau une grandeur semblable à $P$, l'incertitude attendue & Proposé dans un modèle probabiliste de l'attention. Il n'existe pas de mesure directe montrant que le niveau d'acétylcholine suit l'incertitude de l'estimation & \cite{YuDayan2005} & hypothèse \\
13 & Une partie des neurones \nt{ACET} répond à la récompense et à la punition en une vingtaine de millisecondes, d'autant plus fort que le renforcement est inattendu & Souris, neurones cholinergiques du télencéphale basal identifiés par optogénétique: réponse à la récompense et à la punition après $18\pm3$\,ms, amplitude croissant avec le caractère inattendu du renforcement, réponses semblables chez les neurones de deux noyaux & \cite{Hangya2015} & mesuré \\
14 & \nt{NORA} remarque un changement inattendu du monde, \nt{ACET} maintient le réseau en mode écriture & Division du travail proposée dans un modèle. Indirectement: chez l'homme, le diamètre de la pupille, lié à \nt{NORA}, suit les changements et la vitesse d'apprentissage. Le couple \nt{NORA}--\nt{ACET} n'a pas été testé directement & \cite{YuDayan2005,Nassar2012} & hypothèse \\
15 & En sommeil lent, le réseau en mode lecture rejoue ce qui a été écrit pendant la journée, et ces répétitions le recopient dans la mémoire à long terme & Enregistrements d'ensembles de neurones de l'hippocampe du rat: les cellules qui ont travaillé ensemble pendant la journée se déclenchent plus souvent ensemble pendant le sommeil lent qui suit, et dans le même ordre; c'est mesuré. Pour la recopie: supprimer les ondulations hippocampiques (ripples) après l'apprentissage dégrade la mémoire spatiale, et allonger les ondulations spontanées l'améliore; que la cause soit le bas niveau d'\nt{ACET} n'est pas démontré & \cite{WilsonMcNaughton1994,Skaggs1996,Girardeau2009,FernandezRuiz2019,Hasselmo1999} & hypothèse \\
\bottomrule
\end{longtable}}
